\begin{pastexam}[label={px:gw3-100-ds-05}]{關務三等資訊處理 100 年 第五題}
（申論；15 分）

題組：請根據下圖回答第三至五題：

圖中表示有 4 個朋友，名（name）叫張三（Chang San, CS）、李四（Lee Si, LS）、王五（Wang Wu, WW）及趙六（Chao Liu, CL），用兩個字母代表各人，如 CS 代表張三，並標示兩兩住家交通狀況及距離（distance），如張三家有 7 公里的路到李四家。
\begin{center}
\begin{tikzpicture}[x=1cm,y=1cm,v/.style={font=\small\bfseries,inner sep=2pt},>={Stealth[length=2.2mm]},w/.style={font=\small\bfseries,inner sep=2pt}]
\node[v] (CS) at (0,0) {張三};
\node[v] (LS) at (2.4,1.1) {李四};
\node[v] (WW) at (5,0) {王五};
\node[v] (CL) at (2.4,-1.1) {趙六};
\draw[->,thick] (CS)--node[w,above left]{7}(LS);
\draw[->,thick] (LS)--node[w,above right]{21}(WW);
\draw[->,thick] (LS)--node[w,right]{33}(CL);
\draw[->,thick] (CS)--node[w,below left]{20}(CL);
\draw[->,thick] (CL)--node[w,below right]{6}(WW);
\end{tikzpicture}
\end{center}
\begin{quote}
hash function：h(name) = int (1\textsuperscript{st} char of name)*31+int(2\textsuperscript{nd} char of name)\\
hash table size：11 \quad (index 0 – 10)\\
index：h(name) mod 11\\
int( “C” )=67, int( “S” )=83, int( “L” )=76, int( “W” )=87\\
舉張三為例：\\
index( “CS” )=(int( “C” )*31+int( “S” )) mod 11 = (67*31+83) mod 11 = 4
\end{quote}

找出以張三為 root 的 Prim's Minimal Spanning Tree，要畫出 1. tree\quad 2. priority queue 2 個 data structures 的最後結果。（15 分）
\end{pastexam}

\begin{pastsoln}
答案：生成樹的邊是 張三–李四 7、張三–趙六 20、趙六–王五 6，總權重 $33$；priority queue 最後是空的。

\textbf{解讀}。 最小生成樹是對\textbf{無向圖}定義的；Prim 演算法用在有向圖上沒有意義（有向圖對應的是「最小樹形圖」，要用別的演算法）。
這裡照題意把每條路看成雙向，邊為 CS–LS 7、CS–CL 20、LS–WW 21、LS–CL 33、CL–WW 6。 Prim 演算法見資結書第~24.3~節：從根開始，
每次把「一端在樹上、一端不在」的邊中最輕的那條加入。 priority queue 存還不在樹上的點，鍵是它連到樹上的最輕邊的權重。

\textbf{追蹤}（$k/p$ 表示鍵值和連到樹上的哪一點）：
\begin{center}
\small
\begin{tabular}{lcccl}
\toprule
步驟 & LS & CL & WW & priority queue（鍵由小到大） \\
\midrule
加入 CS（根） & $7/$CS & $20/$CS & $\infty$ & (LS, 7), (CL, 20), (WW, $\infty$) \\
加入 LS（邊 CS–LS 7） & & $20/$CS & $21/$LS & (CL, 20), (WW, 21) \\
加入 CL（邊 CS–CL 20） & & & $6/$CL & (WW, 6) \\
加入 WW（邊 CL–WW 6） & & & & 空 \\
\bottomrule
\end{tabular}
\end{center}
加入 LS 後，LS–CL 33 比 CL 目前的鍵 20 大，不更新；LS–WW 21 讓 WW 的鍵從 $\infty$ 降為 21。 加入 CL 後，CL–WW 6 比 21 小，WW 的鍵降為 6、
改連到 CL。

\textbf{最後結果}：
\begin{center}
\begin{tikzpicture}[x=1cm,y=1cm,v/.style={draw,rounded corners=2pt,inner sep=2pt,font=\small},w/.style={font=\small,fill=white,inner sep=1pt}]
\node[v] (CS) at (0,0) {張三 CS};
\node[v] (LS) at (-1.5,-1.3) {李四 LS};
\node[v] (CL) at (1.5,-1.3) {趙六 CL};
\node[v] (WW) at (1.5,-2.6) {王五 WW};
\draw (CS)--node[w]{7}(LS); \draw (CS)--node[w]{20}(CL); \draw (CL)--node[w]{6}(WW);
\end{tikzpicture}
\qquad\qquad
\begin{tabular}[b]{l}
2. priority queue：空 \\[1em]
\end{tabular}
\end{center}
1. tree：以張三為根，李四、趙六是張三的子節點，王五是趙六的子節點；總權重 $7 + 20 + 6 = 33$。 檢查：$4$ 個點、$3$ 條邊、連通；沒選的兩條
LS–WW 21、LS–CL 33 各自和樹上的邊成環，而且都是環上最重的邊，所以這棵樹是最小生成樹。 （若用不支援降低鍵值的佇列，舊的 (WW, 21) 會在 WW 加入
後才被取出丟掉，最後佇列同樣是空的。）
\end{pastsoln}
